\section{Construction hypotheses for the arithmetic comparison}\label{sz:register}
At a fixed admissible datum, Theorem~\ref{sz:conclusion} deduces the signed
inequality from A, B and C. The sharpenings below apply on the constructed
family as $x$ grows with the permissible parameters and scale exponent fixed.

\begin{longtable}{@{}>{\raggedright\arraybackslash}p{.19\linewidth}@{\hspace{.025\linewidth}}p{.755\linewidth}@{}}
\toprule Step & Deductions and remaining requirements\\\midrule\endhead
Local Taylor bound &
Sections~\ref{sec:unitrow-frame}--\ref{sec:other-witnesses} construct the frame and coefficients at split, inert and ramified primes. The local CM identities, the period normalization and the realization by integral test forms enter separately.\\[4pt]
Spreading &
Section~\ref{app:local-frame-recovery} gives a compact-position frame bound. Section~\ref{spread:normalization} computes the scalar in F6. Sections~\ref{moment:rigidity}--\ref{moment:sieve} derive frame separation, a weighted pair count, the primary character formula and the transfer of local errors into both moments. The actual filling still requires the one-term multiplicity, local kernel bounds, orbital integral comparison and compatible measures. Proposition~\ref{prop:weakerprecision} gives an alternative sufficient accuracy.\\[4pt]
Inner forms &
Section~\ref{tr:parabolic-descent} proves local vanishing on a coprime factorization. The global pretrace expansion and the matching of elliptic terms in \S\ref{tr:geometry} remain inputs.\\[4pt]
Twisted lengths &
Sections~\ref{sec:selmer-polynomial-descent}--\ref{sec:terminal-selmer} construct polynomials for the indicated data. Section~\ref{sec:odd-cocycles-reduction} treats the common cohomological integer. Section~\ref{moment:line} verifies the rank-one Selmer conditions for the explicit main line and computes its split vacuum factor. Attaching the actual representations, identifying the local conditions and proving the first-order congruences remain necessary in R.\\[4pt]
Auxiliary forms &
Section~\ref{pathmath:dual} computes the integral pairing correction. The denominator calculation in \S\ref{pathmath:denom} requires the stated kernel and period identities. The required spanning and size of the auxiliary forms remain part of Aux.\\\bottomrule
\end{longtable}
\paragraph{Hypothesis groups on the constructed scales.}
The original register has four groups in $\mathsf A_{\rm loc}$, eight in
F0--F7, and six in I, R, G, F, K, Aux. For the construction at a sufficiently
large cutoff, \S\ref{moment:rigidity} derives F4 from the fitted model and
free principal level. This leaves seventeen other groups to verify on the
selected datum. The groups share inputs, as described below. Several contain substantial written deductions. In
particular the finite character identification in F7 is now
Lemma~\ref{lem:fitted-primary-character}; its inner-form comparison and
formal-degree identification are still required. The simultaneous choice
of datum in Corollary~\ref{cor:construction-grh} applies to all branches.

\subsection*{Where the written proof still depends on construction estimates}
\paragraph{Shortness and integrality.}
The archimedean chain requires the normalized vacuum and period identities,
their growth and expansion bounds, and the integral branching data. At
small nonsplit primes Appendix~\ref{sec:jets} derives a divisibility
estimate from the stated CM and vacuum identities. The marked places and
the level prime also require the transfer, edge and paired-lift inputs
specified in $\mathsf A_{\rm loc}$. These inputs supply both the gain
removed by $c_B$ and the colength of the submodule where the gain holds.

\paragraph{Spreading.}
The restriction to ramified forms uses the loss of at most
$f d_X/(N-1)$ established at the source. The full-source moment still
requires the fitted filling and Bessel block normalizations, the same-link
and unequal-link pair estimates and the remaining
primitive frame estimates of Appendix~\ref{sz:shells}, and the Siegel and
sieve estimates for the open cell. The near-optimal coefficient in F6 is
used in the dimension bound. On the constructed scales the box, separation and orbit statements in F4 follow in the fitted free-level model by \S\ref{moment:rigidity}. The transfer to the good target uses the
degree matching in F7 and the finite character formula of Lemma~\ref{lem:fitted-primary-character} together with Appendix~\ref{sec:trace}.

\paragraph{Determinant height.}
The regular twist uses Galois, Selmer and first-order congruence statements
in R. Appendix~\ref{sec:finite} supplies simultaneous specialization,
the explicit prime sum for the allowance and the passage to test forms
under its stated arithmetic inputs. The compression of the one-lift
operator still requires the normalized period gluing and exchanges in G.
The term $-A_0$ uses the flat value formula, the local curve hypotheses for the sharp witness and
its realization by test forms in F. K supplies the analytic norm formula
and triangular exchanges. Aux supplies the normalized path period, its
surjectivity, and the denominator and size of the grid forms. These are
the inputs to the one-lift height, its telescoping through $m$ lifts and
the extension of the given family at controlled cost.

\paragraph{Three dependencies shared between the branches.}
The short-basis argument is used at two different levels. The source
construction in A and the grid in Aux have different box counts and
require their own conductor and norm estimates. The inner-form comparison
in B at the identity transfers a dimension. In C a determinant transfer
requires the traces of the Hecke operators as well. The common geometric
hypotheses and the extra step for the determinant are in
Appendix~\ref{sec:trace}. Finally, a universal lower bound on the
valuation of a short-form jet in A does not exhibit any coefficient
attaining a small valuation in C. Appendix~\ref{sec:jets} specifies the
additional witness and the arithmetic cases it requires.

The remaining
mathematical work for this exposition lies upstream of \eqref{sz:signed}.
Theorem~\ref{sz:conclusion} applies datum by datum.
For the consequence about Dirichlet zeros, Corollary~\ref{cor:construction-grh}
requires the listed inputs on one constructed datum at each sufficiently
large integer length, for each fixed $(q_0,d,\eta)$.
